\chapter{Duality Defects, Gauging, And Orbifold Data}
\label{ch:global-duality-defects-gauging-orbifold-data}

\ConventionsMixed

Gauging a finite symmetry and placing duality defects are operations on QFT
data.  This chapter treats them as manipulations of topological defects,
background fields, and sums over sectors, with anomaly and boundary choices
made explicit.

\section{Finite Symmetry Backgrounds}

Let \(G\) be a finite internal symmetry of a \(D\)-dimensional QFT
\(\mathcal Q\).  A background \(G\)-field on spacetime \(M\) is a principal
\(G\)-bundle
\[
  P\to M,
\]
or equivalently a Cech one-cocycle \(g_{ij}\) on a good cover, modulo
coboundaries.  The partition function in the background is
\[
  Z_{\mathcal Q}[P].
\]
The anomaly-free gauged theory is defined by the finite sum
\[
  Z_{\mathcal Q/G}(M)
  =
  \frac1{|G|}
  \sum_{[P]\in\operatorname{Bun}_G(M)}
  \frac{Z_{\mathcal Q}[P]}{|\operatorname{Aut}(P)|/|G|},
\]
with the normalization chosen by the convention for the empty gauge field on
connected \(M\).  Equivalent groupoid-sum conventions differ by invertible
topological factors.

\section{Dijkgraaf--Witten Twist}

A class
\[
  \omega\in H^D(BG,U(1))
\]
defines an invertible topological action for the background \(P\).  The
twisted gauging is
\[
  Z_{\mathcal Q/G,\omega}(M)
  =
  \sum_{[P]}\frac1{|\operatorname{Aut}(P)|}
  Z_{\mathcal Q}[P]\,\ee^{\ii S_\omega(P)}.
\]
If \(\mathcal Q\) has anomaly \(\alpha\in H^{D+1}(BG,U(1))\), gauging \(G\)
requires a trivialization of \(\alpha\) or a bulk theory whose boundary is
\(\mathcal Q\).  A Dijkgraaf--Witten twist changes the gauged theory but does
not by itself cancel a nontrivial anomaly one degree higher.

\section{Gauging Interface}

There is a codimension-one interface
\[
  \mathcal N_G:\mathcal Q\longrightarrow \mathcal Q/G
\]
that implements the sum over \(G\)-bundles on one side.  The reverse
interface \(\mathcal N_G^\dagger\) satisfies the fusion relation
\[
  \mathcal N_G^\dagger\circ\mathcal N_G
  =
  \bigoplus_{g\in G}D_g,
\]
where \(D_g\) is the invertible symmetry defect of \(\mathcal Q\).  This
follows by cutting spacetime along two nearby interfaces: the gauge field in
the thin slab is flat, so its holonomy label is an element \(g\in G\), and
each holonomy sector inserts \(D_g\).

The opposite fusion,
\[
  \mathcal N_G\circ\mathcal N_G^\dagger,
\]
is an algebra object in the defect category of the gauged theory.  It
contains the quantum symmetry defects acting on twisted sectors.

\section{Orbifold Local Operators In Two Dimensions}

For a two-dimensional QFT with finite symmetry \(G\), gauging produces local
operators from sectors with monodromy
\[
  g\in G
\]
around the insertion.  If \(hgh^{-1}=g'\), the sectors labelled by \(g\) and
\(g'\) are identified.  The local operator space of the gauged theory has the
form
\[
  \mathcal V_{\mathcal Q/G}
  =
  \bigoplus_{[g]}
  \mathcal V_g^{C_G(g)},
\]
where \(\mathcal V_g\) is the \(g\)-twisted local space and \(C_G(g)\) is the
centralizer.  The operator product is defined by bringing together punctures
with compatible flat \(G\)-bundles on the pair of pants.

\section{Duality Defects}

A duality defect is a topological defect \(D\) such that fusing it with its
orientation reverse produces a sum of symmetry defects rather than the vacuum
defect alone:
\[
  D^\dagger\otimes D\simeq \bigoplus_{g\in G}D_g.
\]
The defect is therefore noninvertible unless \(G\) is trivial.  Its action on
local operators maps untwisted operators into the \(G\)-invariant part of
twisted or gauged sectors, depending on which side of the interface the
operator lies.

\begin{proposition}[Self-dual gauging produces a noninvertible defect]
\label{prop:self-dual-gauging-noninvertible-defect}
Let \(H\) be finite.  Assume that \(H\) acts internally on a QFT
\(\mathcal Q\) and can be gauged without an anomaly.  Suppose there is an
equivalence
\[
  \Phi:\mathcal Q/H\longrightarrow \mathcal Q
\]
after all background choices and local counterterms have been fixed.  Then
the composite
\[
  D_H:=\Phi\circ\mathcal N_H
\]
is a topological defect of \(\mathcal Q\) satisfying
\[
  D_H^\dagger\otimes D_H
  \simeq
  \bigoplus_{h\in H}D_h .
  \label{eq:self-dual-gauging-defect-fusion}
\]
If \(H\) is nontrivial and the \(D_h\) are distinct simple invertible
defects, \(D_H\) is noninvertible.
\end{proposition}

\begin{proof}
The equivalence \(\Phi\) is invertible as an interface, so
\(\Phi^\dagger\circ\Phi\simeq\mathbf 1_{\mathcal Q/H}\).  Therefore
\[
  D_H^\dagger\otimes D_H
  =
  \mathcal N_H^\dagger\circ\Phi^\dagger\circ\Phi\circ\mathcal N_H
  \simeq
  \mathcal N_H^\dagger\circ\mathcal N_H .
\]
The slab calculation of the gauging interface gives
\[
  \mathcal N_H^\dagger\circ\mathcal N_H
  \simeq
  \bigoplus_{h\in H}D_h,
\]
because the flat \(H\)-bundle in the thin slab is classified by its holonomy
\(h\).  If \(D_H\) were invertible, the left side of
\eqref{eq:self-dual-gauging-defect-fusion} would be the simple vacuum defect.
For nontrivial \(H\) with distinct simple \(D_h\), the right side has more
than one simple summand, so no inverse exists.
\end{proof}

\begin{example}[Normal-subgroup gauging interface]
\label{ex:normal-subgroup-gauging-defect}
Let \(G\) be a finite internal symmetry and \(H\triangleleft G\) a normal
subgroup that is anomaly-free as a symmetry of \(\mathcal Q\).  Gauging \(H\)
produces \(\mathcal Q/H\) with residual ordinary symmetry \(G/H\) when no
extension anomaly obstructs the quotient action.  The interface
\(\mathcal N_H:\mathcal Q\to\mathcal Q/H\) has junction operators with the
invertible \(H\)-symmetry defects:
\[
  D_h\otimes\mathcal N_H\simeq\mathcal N_H,
  \qquad h\in H .
\]
At a self-dual point where an equivalence \(\Phi:\mathcal Q/H\to\mathcal Q\)
is part of the QFT datum, Proposition~\ref{prop:self-dual-gauging-noninvertible-defect}
turns the gauging interface into a genuine noninvertible defect of
\(\mathcal Q\).  The critical Ising Kramers--Wannier defect is the case
\(H=\mathbb Z_2\), with fusion
\[
  \mathcal N\otimes\mathcal N\simeq\mathbf 1\oplus\eta .
\]
\end{example}

\section{Gauge-Theory Duality Walls}

In gauge theories, duality walls can act on line-operator lattices.  Let
\(\Gamma\) be the electric-magnetic charge lattice with antisymmetric pairing
\[
  \langle\ ,\ \rangle:\Gamma\times\Gamma\to\mathbb Z.
\]
A duality wall induces an automorphism
\[
  \gamma:\Gamma\to\Gamma,\qquad
  \langle\gamma x,\gamma y\rangle=\langle x,y\rangle .
\]
If the wall is topological, it maps genuine line operators to genuine line
operators after global form and discrete theta data are transformed.  If
additional degrees of freedom live on the wall, their anomaly inflow must
cancel the mismatch between the two sides.

\begin{example}[Four-dimensional abelian line-lattice duality walls]
\label{ex:abelian-line-lattice-duality-walls}
For four-dimensional \(U(1)\) gauge theory with both electric and magnetic
line operators allowed, take
\[
  \Gamma=\mathbb Z^2,\qquad
  x=(q_e,q_m),
\]
with Dirac pairing
\[
  \langle x,y\rangle=q_e q'_m-q_m q'_e .
\]
The two standard generators of \(SL(2,\mathbb Z)\) act by
\[
  S(q_e,q_m)=(q_m,-q_e),
  \qquad
  T(q_e,q_m)=(q_e+q_m,q_m).
\]
Both preserve the pairing.  A topological \(S\)-wall maps an electric Wilson
line to a magnetic 't Hooft line and a magnetic line to the inverse electric
line.  A \(T\)-wall implements the Witten-effect shift of electric charge by
magnetic charge.  These walls are genuine defects only after the global form
of the gauge group, the allowed line lattice, and the discrete theta angle
are transformed consistently; otherwise the wall maps a genuine line on one
side to a line that is not genuine on the other side.
\end{example}

\begin{proposition}[Pairing preservation for the \(S\) and \(T\) walls]
\label{prop:st-walls-preserve-dirac-pairing}
The transformations in Example~\ref{ex:abelian-line-lattice-duality-walls}
preserve the Dirac pairing on \(\Gamma\).
\end{proposition}

\begin{proof}
For \(Sx=(q_m,-q_e)\) and \(Sy=(q'_m,-q'_e)\),
\[
  \langle Sx,Sy\rangle
  =
  q_m(-q'_e)-(-q_e)q'_m
  =
  q_e q'_m-q_m q'_e
  =
  \langle x,y\rangle .
\]
For \(Tx=(q_e+q_m,q_m)\) and \(Ty=(q'_e+q'_m,q'_m)\),
\[
  \langle Tx,Ty\rangle
  =
  (q_e+q_m)q'_m-q_m(q'_e+q'_m)
  =
  q_e q'_m-q_m q'_e
  =
  \langle x,y\rangle .
\]
Thus both maps are automorphisms of the charge lattice with its
antisymmetric Dirac pairing.
\end{proof}

\begin{openproblem}[Continuum gauging interfaces]
\label{op:continuum-gauging-interfaces}
Construct gauging interfaces and duality defects in interacting continuum QFTs
from background-field sums, local operator sectors, defect junctions, anomaly
trivializations, and reflection-positive state spaces.
\end{openproblem}
